\chapter{Card Recognition Pipeline}\label{chap:rec}

\section{Data Acquisition}
One of the primary challenges in training deep learning models for \ac{TCG} identification is \ii{data scarcity}. While the total number of unique cards in circulation is vast (often exceeding tens of thousands for games like \ac{MTG}), labeled datasets containing real-world photographs of these cards are virtually non-existent. Furthermore, collecting and manually labeling a sufficiently large dataset of physical cards is prohibitively expensive and time-consuming.

To address this limitation, this thesis adopts a \ii{synthetic data generation} approach. High-quality reference images, serving as ground truth templates, were retrieved via the Scryfall Bulk Data API~\cite{scryfall}. This service provides a comprehensive catalogue of all printed cards in a structured JSON format, enabling the automated generation of training samples.

\subsection{Scraping Pipeline and Indexing}
Custom data acquisition script, \texttt{data\_scraper.py}, was developed to process the metadata and download high-resolution scans. A critical aspect of the pipeline is the \ii{file indexing strategy}, designed to balance machine uniqueness with human readability, a feature essential for manual inspection and debugging. % TODO: be more descriptive about where the source code is located? scripts/data_scraper.py

Instead of relying solely on human-unreadable identifiers (such as \acp{UUID}), a hybrid naming convetion derived from the card's canonical \ac{URI} is proposed. Each file name combines the card edition, collector number, and name with a \ac{UUID}: 
\begin{equation}
    \texttt{filename} = \texttt{edition\_collectorNumber\_name-UUID.png}
\end{equation}
For example, a card from the \ii{Zendikar} set is stored as:
\begin{center}
    \texttt{zen\_21\_kor-outfitter-00006596-1166-4a79-8443-ca9f82e6db4e.png}
\end{center}
This format allows for immediate visual identification of the card's edition and variant without requiring a lookup table, while the \ac{UUID} suffix guarantees collision-free storage even for edge cases (e.g., duplicate names across sets).


% TODO: asi to chce zmínit, že je důležité nezapomenout trénovat _front a _back zvlášť,
% zároveň mi tohle už nepřijde už akademicky znějící?
\subsection{Handling Complex Layouts}
The script also implements logic to handle complex card layouts, such as \acp{DFC} (transform, modal \acp{DFC}). These entries, identified by the \texttt{layout} property in the JSON metadata, contain multiple faces within a single object. Our pipeline detects these instances and automatically splits them into separate template images, appending \texttt{\_front} or \texttt{\_back} suffixes to their identifiers to ensure that both sides of the card are learnable by the model.

\section{Synthetic Data Generation}
To address the lack of labeled real-world data, the training pipeline incorporates an on-the-fly synthesis module that simulates artifacts typical of physical image acquisition.

For every training iteration, a canonical digital scan is processed through a stochastic augmentation chain. The applied transformations focus on three main categories:
\begin{itemize}
    \item Geometric Transformations: Random perspective warping and rotation are applied to induce invariance to camera pose and alignment errors.
    \item Photometric Distortions: Variations in brightness, contrast, and saturation, combined with Gaussian blur, approximate diverse environmental lighting and optical imperfections.
    \item Occlusions and Artifacts: Random erasure methods (e.g., \texttt{CoarseDropout}~\footnote{Documentation available at \url{https://albumentations.ai/docs/api-reference/albumentations/augmentations/dropout/coarse_dropout/\#CoarseDropout}}) are utilized to simulate physical damage, dust, or partial occlusion of the card surface.
\end{itemize}
Consequently, the model is trained on a virtually infinite stream of unique samples. This strategy minimizes overfitting and enhances the model's generalization capabilities on unseen physical cards.

\section{Theoretical Framework}


\section{Model Architecture}
The implemented recognition pipeline, defined in \texttt{model.py}, consists of three distrinct stages.

\subsection{Backbone Network}
A ResNet-50~\cite{He2015} extractor pretrained on ImageNet.

\subsection{Embedding Head}
A fully connected layer mapping the 2048 backbone features to a compact 512-dimensional embedding space.

\subsection{ArcFace Head}
The \texttt{ArcFaceHead} module, 


During inference (see \texttt{inference.py}), the ArcFace head is discarded, and similarity is computed using the cosine distance between the 512-dimensional embeddings.

